\begin{afterword}
\summaryhead

The post opens with John McCarthy's line ``When there's a will to fail, obstacles can be found.'' The author says he remembered Luke Skywalker from childhood as a hero, and on rewatching found him whiny. He then looked up the scene in \textsc{The Empire Strikes Back} where Yoda says ``Do, or do not. There is no try,'' and Luke, after the X-wing starts to rise from the swamp and sinks back, says it is too big and gives up.

The rest of the post is an invented outtake. The actor Mark Hamill argues with George Lucas that Luke should try again, try for a month, or at least try for five minutes. Lucas refuses each time: he will not halt the story, and the audience will not notice anything odd. His last line gives the moral: people would not try for five minutes before giving up even if the fate of humanity were at stake.

\responsehead

The post is a joke with a thesis in its last line, and neither holds up.

The joke misreports its target. In the scene it mocks, Luke's giving up is not presented as ordinary. Yoda immediately lifts the ship out of the swamp, Luke says he does not believe it, and Yoda answers, ``That is why you fail.'' The film shows the surrender, calls it failure, and names a cause. The outtake's complaints (Luke will be stranded, the audience will not buy it) depend on leaving that out. The film's diagnosis is also not the post's. Yoda's lesson is about belief, not stamina. The post turns a scene about faith into a scene about effort and then faults it for not teaching effort. Its one attribution is wrong as well: Lucas did not direct the film, and the scene is not his to defend.

The thesis is generalization from fictional evidence. ``People wouldn't try for five minutes before giving up'' is a claim about human beings. The only support offered is a paced film scene, and the objection is built into the parody: the Lucas character says he will not halt the story. The author had named this exact fallacy a year earlier, and said then that the needs of storytelling are not those of forecasting. Here he uses a story's pacing as a fact about people and puts the conclusion in the mouth of a real man who never said it.

Then there is the move in the last clause. The scene concerns one man and his spaceship. The moral concerns ``the fate of humanity.'' The jump carries in the author's larger message: that the reader who stops trying is failing to save the world. That claim is not argued anywhere in the post. It arrives as a punch line, and the reader is invited to hold ``a pathetic loser'' in contempt on the way.

The context makes the post weaker still. Its argument lives in ``Trying to Try,'' published the day before. The Highlights collection includes this post and omits that one. What the collection's readers get is a fan's complaint and a claim about human nature, with the reasoning removed.

In short: a parody that misreports the scene it mocks, and a claim about human nature drawn from a film, the fallacy its author had named a year before.
\end{afterword}
